\section*{Legacy problem audit for VI/15}
\addcontentsline{toc}{section}{Legacy problem audit for VI/15}

The Volume VI problem ledger was consulted before choosing any canonical Problem count.  The
coverage inventory identifies exactly three problem/exercise/solution source scopes for this
chapter: the preferred reorganized sheaf notes, AG1's kernel/image/quotient selector, and
AG2's kernel/image/quotient selector.

\begin{description}[style=nextline]
\item[AG1 Exercise 5 --- INCLUDE/MERGE.]
Its mathematical center is the image sheaf as the subsheaf of locally liftable sections.
It becomes VI.15.P1 as a full detailed dossier.  The parallel proof in the preferred
reorganized notes is merged into the chapter and P1 rather than duplicated.

\item[AG2 Problem 8 --- INCLUDE.]
Its mathematical center is the quotient sheaf and the relationship with short exact
quotient sequences.  It becomes VI.15.P2.  Both original directions are retained: quotient
by a subsheaf gives the quotient sequence, and conversely the final term of a short exact
sequence is the quotient by the image of the first term.

\item[Preferred sheaf notes T04 --- MERGE.]
The kernel, image-presheaf versus image-sheaf, quotient, and cokernel constructions are
migrated into the chapter theory, example bank, VI.15.P1, and VI.15.P2.  They do not create
a second copy of either numbered legacy problem.

\item[AG1 Exercise 6 and AG2 Problem 7 --- MOVE.]
These concern the general stalkwise exactness/isomorphism theory and remain reserved for
VI/16.  They are not consumed by the quotient-sequence preview required by AG2 Problem 8.

\item[AG2 Problem 9 --- MOVE.]
The left exactness of $\Gamma(U,-)$ remains reserved for VI/48 in the cohomology arc.

\item[Merged-exercise archive lineage --- ARCHIVE-SOURCE-ONLY.]
The repository duplicate map classifies the historical merged/corrected exercise snapshots
as source-only lineage.  The already-audited representative variants do not expose a unique
sheaf-family Problem requiring a second VI/15 destination; provenance remains preserved in
the archive lineage.
\end{description}

No substantive numbered corpus problem assigned to VI/15 is discarded.  The canonical
count is therefore corpus-driven plus new synthesis: two preserved legacy problems and
eighteen new synthesis dossiers, for twenty Problems in this chapter.  The global ledger is
updated so AG1 Exercise 5 and AG2 Problem 8 are no longer merely reserved; they now have
canonical destinations VI.15.P1 and VI.15.P2.
